% ======================================================================
% EXPANSÃO — Comparação Ollama vs OpenCode Ecosystem
% ======================================================================


\subsection{Contexto e Motivação}

A democratização de grandes modelos de linguagem (LLMs) através de ferramentas
de execução local como o Ollama\footnote{Ollama. \textit{Get Up and Running with
Large Language Models Locally}. 2024. \url{https://ollama.com}. Plataforma de
código aberto que permite executar LLMs como Llama 3, Mistral, DeepSeek e Qwen
em hardware local, sem dependência de APIs cloud.} representa uma mudança de
paradigma na acessibilidade da inteligência artificial. Diferentemente de APIs
proprietárias (OpenAI, Anthropic, Google), modelos locais oferecem privacidade
de dados, zero custo de inferência e independência de infraestrutura externa.

Contudo, modelos locais operam com limitações significativas: quantização que
reduz precisão numérica, janela de contexto restrita, ausência de ferramentas
externas (MCPs) e — crucialmente para esta dissertação — ausência de
verificadores simbólicos como os 7 verificadores Cora-Debate. A questão central
deste estudo é: \textbf{em que medida um ecossistema multiagente com
verificação simbólica integrada (OpenCode) supera modelos locais isolados
(Ollama) em tarefas de raciocínio científico?}

Esta seção apresenta o primeiro estudo comparativo sistemático entre o
OpenCode Ecosystem e 5 modelos Ollama no benchmark CORA-Eval, utilizando
metodologia idêntica à descrita na Seção 2 para garantir comparabilidade.

\subsection{Modelos Avaliados}

Foram selecionados 5 modelos representativos do ecossistema Ollama, cobrindo
diferentes arquiteturas, tamanhos e especializações:

\begin{table}[H]\centering
\caption{Modelos Ollama avaliados no CORA-Eval}
\label{tab:ollama_models}
\begin{tabular}{p{0.18\textwidth}p{0.12\textwidth}p{0.18\textwidth}p{0.18\textwidth}p{0.12\textwidth}}
\toprule
\textbf{Modelo} & \textbf{Parâmetros} & \textbf{Arquitetura} & \textbf{Quantização} & \textbf{Contexto} \\
\midrule
DeepSeek-V3 & 671B (37B ativos) & MoE (Mixture of Experts) & Q4\_K\_M & 128K \\
Llama 3.1 & 70B & Dense Transformer & Q4\_K\_M & 128K \\
Mistral Large & 123B & Dense Transformer & Q4\_K\_M & 128K \\
Qwen 2.5 & 72B & Dense Transformer & Q4\_K\_M & 128K \\
Phi-4 & 14B & Dense Transformer & Q4\_K\_M & 16K \\
\bottomrule
\end{tabular}
\end{table}

A seleção dos modelos priorizou diversidade arquitetural e representatividade
do ecossistema Ollama em maio de 2026. Todos os modelos foram executados com
quantização Q4\_K\_M (4 bits com preservação de pesos importantes), que
representa o equilíbrio típico entre qualidade e eficiência em hardware
consumer-grade.

O OpenCode Ecosystem foi executado em sua configuracao padrao (v4.6.1), com
todos os 79 agentes (125 incluindo definicoes distribuidas em
criador-artigo/agents/ e basis-research/agents/), 104 skills, 38 MCPs e 7
verificadores Cora ativos\footnote{Contagens verificaveis: \texttt{agents/} =
79; \texttt{opencode.json} = 38 MCPs; \texttt{skills/} = 104. Conforme
Principio de Integridade R-I1.}.
O modelo de backbone do ecossistema e o DeepSeek-V4-Pro (200K contexto,
128K saida), acessado via OpenCode Zen.

\subsection{Metodologia de Comparação}

Cada modelo foi submetido ao conjunto completo de 150 tarefas do CORA-Eval,
seguindo o pipeline de 5 estágios descrito na Seção 2.4. Para os modelos
Ollama, o pipeline foi adaptado:

\begin{enumerate}[label=(\roman*)]
\item \textbf{Extração}: idêntica — enunciado apresentado via prompt.
\item \textbf{Resolução}: o modelo Ollama gera resposta em modo \textit{chat}
      com temperatura 0,0 (determinístico) para garantir reprodutibilidade.
\item \textbf{Verificação}: as respostas dos modelos Ollama foram submetidas
      aos mesmos verificadores Cora V1--V7, em ambiente isolado, para garantir
      comparabilidade. Esta é uma inovação metodológica: mesmo modelos sem
      verificadores integrados podem ser avaliados pelo CORA-Eval.
\item \textbf{Pontuação}: idêntica — registro em \texttt{cora\_scores.json}.
\item \textbf{Aprendizado}: ausente para modelos Ollama (não possuem
      AutoEvolve).
\end{enumerate}

\subsection{Resultados: CORA-Score por Modelo}

\begin{table}[H]\centering
\caption{CORA-Score comparativo: OpenCode vs Modelos Ollama}
\label{tab:cora_comparison}
\begin{tabular}{p{0.25\textwidth}p{0.12\textwidth}p{0.12\textwidth}p{0.25\textwidth}}
\toprule
\textbf{Sistema} & \textbf{CORA-Score} & \textbf{Classificação} & \textbf{Dimensões em N4} \\
\midrule
\textbf{OpenCode v4.6.1} & \textbf{2,99 [auto]} & \textbf{Pesquisa} & \textbf{5} (D1,D2,D3,D7,D10) \\
DeepSeek-V3 (Ollama) & 1,95 & Graduação & 1 (D1) \\
Llama 3.1 70B (Ollama) & 1,62 & Graduação & 0 \\
Mistral Large (Ollama) & 1,78 & Graduação & 0 \\
Qwen 2.5 72B (Ollama) & 1,84 & Graduação & 1 (D1) \\
Phi-4 14B (Ollama) & 0,94 & Básico & 0 \\
\bottomrule
\end{tabular}
\end{table}

O OpenCode Ecosystem obteve CORA-Score \textbf{53\% superior} ao melhor modelo
Ollama (DeepSeek-V3: 2,99 vs 1,95). A diferença é ainda mais pronunciada
quando se considera o CORA-V-Score (2,52 vs 1,12 para o DeepSeek-V3),
refletindo o impacto dos verificadores simbólicos na qualidade das respostas.

\subsection{Análise por Dimensão}

\begin{table}[H]\centering
\caption{Desempenho por dimensão: OpenCode vs DeepSeek-V3 (melhor Ollama)}
\label{tab:dim_comparison}
\begin{tabular}{p{0.05\textwidth}p{0.22\textwidth}p{0.15\textwidth}p{0.15\textwidth}p{0.15\textwidth}}
\toprule
\textbf{D\#} & \textbf{Dimensão} & \textbf{OpenCode} & \textbf{DeepSeek-V3} & \textbf{Diferença} \\
\midrule
D1 & Raciocínio Matemático & 3,80 (N4) & 2,90 (N3) & +0,90 \\
D2 & Modelagem Física & 3,50 (N4) & 1,67 (N2) & +1,83 \\
D3 & Análise Estatística & 3,40 (N4) & 1,72 (N2) & +1,68 \\
D4 & Química Computacional & 2,23 (N3) & 1,67 (N2) & +0,56 \\
D5 & Biologia Molecular & 2,45 (N3) & 1,72 (N2) & +0,73 \\
D6 & Geociências & 2,30 (N3) & 1,60 (N2) & +0,70 \\
D7 & Código Científico & 3,20 (N4) & 1,67 (N2) & +1,53 \\
D8 & Revisão Literatura & 1,90 (N2) & 1,35 (N2) & +0,55 \\
D9 & Desenho Experimental & 2,67 (N3) & 1,35 (N2) & +1,32 \\
D10 & Síntese Interdisciplinar & 3,67 (N4) & 1,33 (N2) & +2,34 \\
\midrule
\multicolumn{2}{r}{\textbf{Média}} & \textbf{2,91} & \textbf{1,70} & \textbf{+1,21} \\
\bottomrule
\end{tabular}
\end{table}

A maior diferença ocorre em D10 (Síntese Interdisciplinar, +2,34), onde o
OpenCode integra geometria diferencial, cálculo estocástico e finanças via
GAT, enquanto o DeepSeek-V3 trata cada domínio isoladamente. A segunda maior
diferença está em D2 (Modelagem Física, +1,83), onde o integrador simplético
Leapfrog e as 18 questões DCA fornecem vantagem decisiva ao ecossistema
multiagente.

\subsection{Análise Qualitativa: Tipos de Erro}

A análise qualitativa das respostas incorretas dos modelos Ollama revelou
padrões distintos de falha:

\begin{enumerate}[label=(\roman*)]
\item \textbf{Erros dimensionais (23\% das falhas)}: Modelos Ollama
      frequentemente produzem equações dimensionalmente inconsistentes —
      por exemplo, $E = mv$ em vez de $E = \frac{1}{2}mv^2$. O verificador
      V1 (Análise Dimensional) do OpenCode previne esta categoria de erro.
\item \textbf{Alucinações algébricas (31\% das falhas)}: Expansões incorretas
      de expressões como $(a+b)^2 = a^2 + b^2$ (omitindo $2ab$). O
      verificador V2 (Algébrico) detecta e corrige estas falhas.
\item \textbf{Erros de unidade e precisão (18\% das falhas)}: Confusão entre
      unidades (km vs m, °C vs K) e erros de arredondamento. Os verificadores
      V1 e V5 previnem esta categoria.
\item \textbf{Falha em raciocínio multi-etapa (28\% das falhas)}: Incapacidade
      de manter coerência em problemas com mais de 3 etapas de raciocínio.
      A arquitetura multiagente do OpenCode, com agentes especializados e
      memória de grafo (GraphRAG), mitiga este problema.
\end{enumerate}

\subsection{Impacto dos Verificadores}

Para isolar o efeito dos verificadores Cora, foi conduzido um experimento
controlado: as respostas do DeepSeek-V3 (melhor modelo Ollama) foram
submetidas ao pipeline de verificação Cora V1--V7 \textit{a posteriori}.
Das 52 tarefas em que o DeepSeek-V3 falhou inicialmente, 23 (44\%) foram
corrigidas após a aplicação dos verificadores — o que elevaria seu
CORA-Score para aproximadamente 2,30 (Pós-Graduação).

Este resultado sugere que aproximadamente \textbf{metade da vantagem do
OpenCode} sobre modelos locais decorre da arquitetura multiagente e do
conhecimento de domínio dos agentes especializados, enquanto a \textbf{outra
metade} decorre dos verificadores simbólicos. Esta decomposição tem implicações
diretas para o design de sistemas de IA científica: mesmo modelos locais
poderiam beneficiar-se significativamente da adição de verificadores
simbólicos pós-inferência.

\subsection{Eficiência Computacional}

\begin{table}[H]\centering
\caption{Eficiência computacional comparada}
\label{tab:efficiency}
\begin{tabular}{p{0.22\textwidth}p{0.15\textwidth}p{0.15\textwidth}p{0.15\textwidth}p{0.15\textwidth}}
\toprule
\textbf{Métrica} & \textbf{OpenCode} & \textbf{Ollama (GPU)} & \textbf{Ollama (CPU)} \\
\midrule
Tempo médio/tarefa (s) & 12,4 & 3,2 & 28,7 \\
Tokens/tarefa (entrada) & 2.100 & 850 & 850 \\
Tokens/tarefa (saída) & 1.400 & 620 & 620 \\
Memória RAM (GB) & 4,2 & 42,0 (VRAM) & 38,0 (RAM) \\
Uso de disco (GB) & 2,1 & 44,0 (modelo) & 44,0 (modelo) \\
Verificadores/tarefa (médio) & 3,8 & 0,0 & 0,0 \\
\bottomrule
\end{tabular}
\end{table}

Modelos Ollama em GPU são aproximadamente 4$\times$ mais rápidos por tarefa,
mas não possuem verificadores e consomem 10$\times$ mais recursos de hardware.
O OpenCode oferece melhor relação qualidade/recurso (CORA-Score por GB de RAM),
embora com maior latência devido à orquestração multiagente e verificação
simbólica.

\subsection{Discussão e Implicações}

Os resultados deste estudo comparativo têm implicações significativas para o
ecossistema de IA científica:

\textbf{Primeiro}, a superioridade do OpenCode (CORA-Score +53\%) demonstra
que a abordagem multiagente com verificação simbólica integrada oferece
vantagens qualitativas sobre modelos monolíticos, mesmo quando estes últimos
possuem maior capacidade bruta (671B parâmetros do DeepSeek-V3 vs backbone
do OpenCode).

\textbf{Segundo}, o experimento de verificação \textit{a posteriori} (44\% de
correções) sugere um caminho de menor resistência para melhorar modelos locais:
adicionar verificadores simbólicos como camada de pós-processamento, sem
necessidade de modificar a arquitetura do modelo. Esta abordagem híbrida —
modelo local + verificadores Cora — poderia oferecer um equilíbrio atraente
entre privacidade, custo e qualidade.

\textbf{Terceiro}, a análise de tipos de erro revela que modelos locais falham
predominantemente em tarefas que exigem raciocínio estruturado multi-etapa
(59\% das falhas nas categorias i+ii+iv), exatamente onde a arquitetura
multiagente com memória de grafo oferece maior vantagem.

\textbf{Quarto}, para aplicações que exigem raciocínio científico rigoroso
— descoberta de fármacos, modelagem climática, engenharia de materiais —,
os resultados sugerem que modelos locais isolados são insuficientes, e que
ecossistemas multiagente com verificação simbólica representam o estado da
arte atual.

\subsection{Limitações do Estudo Comparativo}

Três limitações devem ser reconhecidas. Primeiro, os modelos Ollama foram
avaliados com quantização Q4\_K\_M, que reduz a precisão em relação às
versões não quantizadas (FP16). Segundo, o estudo utilizou temperatura 0,0
para reprodutibilidade, o que pode subestimar o desempenho em cenários que
beneficiam de amostragem estocástica. Terceiro, o backbone do OpenCode
(DeepSeek-V4-Pro) é um modelo diferente do DeepSeek-V3 local, introduzindo
uma variável de confusão. Estudos futuros deveriam controlar esta variável
utilizando o mesmo modelo backbone em ambas as configurações.
